% Part: lambda-calculus
% Chapter: syntax
% Section: term-revisited

\documentclass[../../../include/open-logic-section]{subfiles}

\begin{document}

\olfileid{lam}{syn}{tr}
\olsection{ترم‌ها به‌مثابهٔ کلاس‌های هم‌ارزی~$\alpha$}

از این پس ترم‌ها را تا حد هم‌ارزی~$\alpha$ در نظر خواهیم گرفت. یعنی
وقتی ترمی می‌نویسیم، مقصودمان کلاس هم‌ارزی~$\alpha$ای است که آن ترم
در آن قرار دارد. برای مثال، $\lambd[a][\lambd[b][a c]]$ را برای مجموعهٔ
همهٔ ترم‌های هم‌ارز به~$\alpha$ با آن می‌نویسیم، مانند
$\lambd[a][\lambd[b][a c]]$، $\lambd[b][\lambd[a][b c]]$ و جز آن.

همچنین، درحالی‌که در بخش‌های پیشین از حروفی چون~$N, Q$ برای نمایش یک
ترم استفاده می‌شد، از این پس آن‌ها را برای نمایش یک کلاس به کار می‌بریم؛
و در ادامه همین کلاس‌ها، نه ترم‌ها، موضوع بررسی ما خواهند بود. حروفی
چون~$x, y$ همچنان یک متغیر را نشان می‌دهند.

همچنین نماد~$\rep{M}$ را برای نمایش یک !!{element} دلخواه از کلاس~$M$
اختیار می‌کنیم، و اگر به بیش از یکی نیاز داشته باشیم،
$\rep{M}[0], \rep{M}[1], etc. $ را به کار می‌بریم.

برای ساده‌کردن بیان، نمادگذاری ترم‌ها را دوباره به کار می‌گیریم. تعریف
زیر را روی کلاس‌ها داریم:
\begin{defn}
  \begin{enumerate}
  \item $\lambd[x][N]$ به‌صورت کلاسی تعریف می‌شود که
    $\lambd[x][\rep{N}]$ را در بر دارد.
  \item $PQ$ کلاسی تعریف می‌شود که~$\rep{P}\rep{Q}$ را در بر دارد.
  \end{enumerate}
\end{defn}

به‌آسانی می‌توان دید که این تعریف‌ها خوش‌تعریف‌اند، زیرا تبدیل~$\alpha$
سازگار است.

\begin{defn} \ollabel{def:fv}
  \emph{متغیرهای آزاد} یک کلاس هم‌ارزی~$\alpha$ مانند~$M$، یا
  $\FV{M}$، به‌صورت~$\FV{\rep{M}}$ تعریف می‌شود.
\end{defn}

این تعریف خوش‌تعریف است، زیرا چنان‌که در~\olref[alp]{thm:fv} نشان داده
شد، $\FV{\rep{M}[0]} = \FV{\rep{M}[1]}$.

نمادگذاری جای‌گذاری در کلاس‌ها را نیز دوباره به کار می‌گیریم:
\begin{defn} \ollabel{def:sub}
  \emph{جای‌گذاری}~$R$ به‌جای~$y$ در~$M$، یا~$\Subst{M}{R}{y}$،
  به‌صورت کلاس هم‌ارزی آلفا شامل
  $\Subst{\rep{M}}{\rep{R}}{y}$ تعریف می‌شود، برای هر~$\rep{M}$ و
  $\rep{R}$ای که این جای‌گذاری را تعریف‌شده سازند.
\end{defn}

چنان‌که در~\olref[alp]{cor:sub} نشان داده شد، این تعریف نیز خوش‌تعریف است.

توجه کنید که این تعریف چگونه استدلال ما را به‌طور چشمگیری ساده می‌کند.
برای مثال:
\begin{align}
  \Subst{\lambd[x][x]}{y}{x} & =\ollabel{eq:1}\\
  &= \Subst{\lambd[z][z]}{y}{x} \ollabel{eq:2}\\
  &= \lambd[z][\Subst{z}{y}{x}] \\
  &= \lambd[z][z]
\end{align}

اگر همچنان آن را جای‌گذاری روی ترم‌ها بدانیم، \olref{eq:1} تعریف‌نشده
است؛ اما چنان‌که پیش‌تر گفته شد، اکنون آن را جای‌گذاری روی کلاس‌ها در
نظر می‌گیریم و ازاین‌رو~\olref{eq:2} ممکن است: می‌توانیم
$\lambd[x][x]$ را با~$\lambd[z][z]$ جایگزین کنیم، زیرا به یک کلاس تعلق
دارند.

به همین دلیل، از این پس فرض می‌کنیم نماینده‌هایی که برمی‌گزینیم همواره
شرایط لازم برای جای‌گذاری را برآورده می‌کنند. برای مثال، وقتی
$\Subst{\lambd[x][N]}{R}{y}$ را می‌بینیم، فرض می‌کنیم نمایندهٔ
$\lambd[x][N]$ چنان برگزیده شده است که~$x \neq y$ و~$x \notin
\FV{R}$.

چون «کلاس» نامیدن~$\lambd[x][x]$ اندکی نامأنوس است، از این پس آن‌ها را
برای تمایز با ترم‌های~$\lambda$ آشنایمان، ترم‌های~$\Lambda$ (یا در
ادامهٔ این بخش صرفاً «ترم») می‌نامیم.

\begin{editorial}
  هنوز نمی‌توانیم با ترم‌ها وداع کنیم: سراسر تعریف ترم‌های~$\Lambda$
  بر ترم‌های~$\lambda$ استوار است، و روشی برای تعریف توابع روی
  ترم‌های~$\Lambda$ عرضه نکرده‌ایم؛ یعنی همهٔ این توابع باید نخست روی
  ترم‌های~$\lambda$ تعریف شوند و سپس، همان‌گونه که برای جای‌گذاری‌ها
  انجام دادیم، روی ترم‌های~$\Lambda$ «تصویر» شوند. بااین‌حال فرض می‌کنیم
  خواننده می‌تواند به‌طور شهودی دریابد که چگونه می‌توان توابع را روی
  ترم‌های~$\Lambda$ تعریف کرد.
\end{editorial}
\end{document}
